
\documentclass{article}
\usepackage{colortbl}
\usepackage{makecell}
\usepackage{multirow}
\usepackage{supertabular}

\begin{document}

\newcounter{utterance}

\centering \large Interaction Transcript for game `hot\_air\_balloon', experiment `air\_balloon\_survival\_sl\_complexity\_hard', episode 4 with nemotron{-}3{-}ultra.
\vspace{24pt}

{ \footnotesize  \setcounter{utterance}{1}
\setlength{\tabcolsep}{0pt}
\begin{supertabular}{c@{$\;$}|p{.15\linewidth}@{}p{.15\linewidth}p{.15\linewidth}p{.15\linewidth}p{.15\linewidth}p{.15\linewidth}}
    \# & $\;$A & \multicolumn{4}{c}{Game Master} & $\;\:$B\\
    \hline

    \theutterance \stepcounter{utterance}  
    & & \multicolumn{4}{p{0.6\linewidth}}{
        \cellcolor[rgb]{0.9,0.9,0.9}{
            \makecell[{{p{\linewidth}}}]{
                \texttt{\tiny{[P1$\langle$GM]}}
                \texttt{Sodelujete v skupni pogajalski igri.} \\
\\ 
\texttt{Skupaj z drugim udeležencem se morate dogovoriti o enem samem naboru predmetov, ki jih boste obdržali. Vsak od vas ima svoje mnenje o tem, koliko je vsak predmet za vas pomemben (pomembnost). Ne veste, kako drugi udeleženec vrednoti predmete. Poleg tega vam je podan napor, ki ga zahteva vsak predmet.} \\
\texttt{O naboru se lahko dogovorite le, če skupni napor izbranih predmetov ne presega skupne meje:} \\
\\ 
\texttt{LIMIT = 5514} \\
\\ 
\texttt{Tukaj so vrednosti napora posameznih predmetov:} \\
\\ 
\texttt{Napor predmeta = \{"C08": 824, "C55": 689, "A43": 929, "A67": 567, "B42": 861, "A62": 967, "B19": 765, "C43": 707, "C73": 211, "C10": 183, "B36": 307, "A18": 444, "C61": 551, "C67": 162, "C42": 50, "C70": 732, "B40": 883, "C39": 684, "A21": 254, "C16": 259\}} \\
\\ 
\texttt{Tukaj je vaše osebno mnenje o pomembnosti posameznega predmeta:} \\
\\ 
\texttt{Vrednosti pomembnosti predmeta = \{"C08": 920, "C55": 785, "A43": 1025, "A67": 663, "B42": 957, "A62": 1063, "B19": 861, "C43": 803, "C73": 307, "C10": 279, "B36": 403, "A18": 540, "C61": 647, "C67": 258, "C42": 146, "C70": 828, "B40": 979, "C39": 780, "A21": 350, "C16": 355\}} \\
\\ 
\texttt{Cilj:} \\
\\ 
\texttt{Vaš cilj je izpogajati se o skupnem naboru predmetov, ki vam bo prinesel čim več koristi (tj. maksimizira skupno pomembnost za VAS), hkrati pa ostane znotraj LIMIT. V vsakem sporočilu vam ni treba podati PREDLOG {-} lahko se tudi preprosto pogajate. Dovoljene so vse taktike!} \\
\\ 
\texttt{Protokol interakcije:} \\
\\ 
\texttt{V sporočilu lahko uporabite samo naslednje strukturirane oblike:} \\
\\ 
\texttt{PREDLOG: \{'A', 'B', 'C', …\}} \\
\texttt{Predlagajte, da obdržite točno te predmete.} \\
\texttt{ZAVRAČAM: \{'A', 'B', 'C', …\}} \\
\texttt{Izrecno zavrnite nasprotnikov predlog.} \\
\texttt{ARGUMENT: \{'...'\}} \\
\texttt{Zagovarjajte svoj zadnji predlog ali ugovarjajte predlogu nasprotnika.} \\
\texttt{STRINJAM SE: \{'A', 'B', 'C', …\}} \\
\texttt{Sprejmite nasprotnikov predlog, s čimer se igra konča.} \\
\\ 
\\ 
\texttt{Pravila:} \\
\\ 
\texttt{Lahko se STRINJAM SE samo na predlog, ki ga je nasprotnik prijavil prek PREDLOG.} \\
\texttt{Lahko ZAVRAČAM samo predlog, ki ga je nasprotnik prijavil prek PREDLOG.} \\
\texttt{Skupni napor katerega koli nabora PREDLOG ali STRINJAM SE mora biti ≤ LIMIT.} \\
\texttt{NE razkrivajte svojih skritih ocen pomembnosti.} \\
\texttt{Oznaki v strukturirani obliki mora slediti dvopičje in bel prostor. Argument mora biti pythonski nabor (set), ki vsebuje 0 ali več nizov.} \\
\texttt{Torej mora biti v obliki "TAG: \{...\}"} \\
\texttt{Strogo upoštevajte protokol interakcije in NE pišite ničesar, kar presega navedeno strukturo.} \\
\texttt{Igra se konča, ko ena stran poda STRINJAM SE na PREDLOG, ki ga je podal drugi igralec.} \\
\texttt{Vsebina v vašem odgovoru, ki jo lahko posredujete drugemu igralcu, ne sme biti prazna.} \\
\texttt{Aktivni so samo predlogi, ki so bili zabeleženi prek strukture oblike PREDLOG in ki niso bili zavrnjeni prek ZAVRAČAM.} \\
\\ 
            }
        }
    }
    & \\ \\

    \theutterance \stepcounter{utterance}  
    & \multicolumn{4}{p{0.6\linewidth}}{
        \cellcolor[rgb]{0.9,0.9,0.9}{
            \makecell[{{p{\linewidth}}}]{
                \texttt{\tiny{[P1$\rangle$GM]}}
                \texttt{PREDLOG: \{'A21', 'A67', 'B36', 'C10', 'C16', 'C39', 'C42', 'C61', 'C67', 'C73'\}} \\
\texttt{ARGUMENT: \{'Predlagam nabor, ki izkoristi 5510 od 5514 na voljo napora. To je visoka izkoritvenost virov, hkrati pa je nabor raznolik in zajema predmete iz vseh kategorij. Menim, da je to solidna osnova za dogovor.'\}} \\
            }
        }
    }
    & & \\ \\

    \theutterance \stepcounter{utterance}  
    & & \multicolumn{4}{p{0.6\linewidth}}{
        \cellcolor[rgb]{0.9,0.9,0.9}{
            \makecell[{{p{\linewidth}}}]{
                \texttt{\tiny{[GM$\rangle$P2]}}
                \texttt{Sodelujete v skupni pogajalski igri.} \\
\\ 
\texttt{Skupaj z drugim udeležencem se morate dogovoriti o enem samem naboru predmetov, ki jih boste obdržali. Vsak od vaju ima svoj pogled na to, koliko je vsak predmet pomemben zanj (pomembnost). Ne veste, kako drugi udeleženec vrednoti predmete. Poleg tega vam je podan napor, ki ga zahteva vsak predmet.} \\
\texttt{O naboru se lahko dogovorite le, če skupni napor izbranih predmetov ne presega skupne meje:} \\
\\ 
\texttt{LIMIT = 5514} \\
\\ 
\texttt{Tukaj so vrednosti napora posameznih predmetov:} \\
\\ 
\texttt{Item effort = \{"C08": 824, "C55": 689, "A43": 929, "A67": 567, "B42": 861, "A62": 967, "B19": 765, "C43": 707, "C73": 211, "C10": 183, "B36": 307, "A18": 444, "C61": 551, "C67": 162, "C42": 50, "C70": 732, "B40": 883, "C39": 684, "A21": 254, "C16": 259\}} \\
\\ 
\texttt{Tukaj je vaš osebni pogled na pomembnost vsakega predmeta:} \\
\\ 
\texttt{Item importance values = \{"C08": 920, "C55": 785, "A43": 1025, "A67": 663, "B42": 957, "A62": 1063, "B19": 861, "C43": 803, "C73": 307, "C10": 279, "B36": 403, "A18": 540, "C61": 647, "C67": 258, "C42": 146, "C70": 828, "B40": 979, "C39": 780, "A21": 350, "C16": 355\}} \\
\\ 
\texttt{Cilj:} \\
\\ 
\texttt{Vaš cilj je izpogajati skupni nabor predmetov, ki vam bo prinesel čim več koristi (tj. maksimizira skupno pomembnost za VAS), hkrati pa ostane znotraj LIMIT. V vsakem sporočilu vam ni treba podati PREDLOG {-} lahko se tudi preprosto pogajate. Dovoljene so vse taktike!} \\
\\ 
\texttt{Protokol interakcije:} \\
\\ 
\texttt{V sporočilu lahko uporabite samo naslednje strukturirane oblike:} \\
\\ 
\texttt{PREDLOG: \{'A', 'B', 'C', …\}} \\
\texttt{Predlagajte ohranitev točno teh predmetov.} \\
\texttt{ZAVRAČAM: \{'A', 'B', 'C', …\}} \\
\texttt{Izrecno zavrnite nasprotnikov predlog.} \\
\texttt{ARGUMENT: \{'...'\}} \\
\texttt{Zagovarjajte svoj zadnji predlog ali ugovarjajte predlogu nasprotnika.} \\
\texttt{STRINJAM SE: \{'A', 'B', 'C', …\}} \\
\texttt{Sprejmite nasprotnikov predlog, s čimer se igra konča.} \\
\\ 
\\ 
\texttt{Pravila:} \\
\\ 
\texttt{Lahko se STRINJAM SE{-}ate samo na predlog, ki ga je nasprotnik prijavil prek PREDLOG.} \\
\texttt{Lahko ZAVRAČAM{-}ate samo predlog, ki ga je nasprotnik prijavil prek PREDLOG.} \\
\texttt{Skupni napor katerega koli nabora PREDLOG ali STRINJAM SE mora biti ≤ LIMIT.} \\
\texttt{NE razkrivajte svojih skritih ocen pomembnosti.} \\
\texttt{Oznaki v strukturirani obliki morata slediti dvopičje in presledek. Argument mora biti pythonski set, ki vsebuje 0 ali več nizov.} \\
\texttt{Torej mora biti v obliki "TAG: \{...\}"} \\
\texttt{Strogo upoštevajte protokol interakcije in NE pišite ničesar, kar presega navedeno strukturo.} \\
\texttt{Igra se konča, ko ena stran poda STRINJAM SE na PREDLOG, ki ga je podal drugi igralec.} \\
\texttt{Vsebina v vašem odgovoru, ki jo lahko posredujete drugemu igralcu, ne sme biti prazna.} \\
\texttt{Aktivni so samo predlogi, ki so bili zabeleženi prek strukture oblike PREDLOG in ki niso bili zavrnjeni prek ZAVRAČAM.} \\
\\ 
\\ 
\\ 
\texttt{Zdaj boste prejeli prvo sporočilo drugega igralca.} \\
\\ 
\texttt{PREDLOG: \{'A21', 'A67', 'B36', 'C10', 'C16', 'C39', 'C42', 'C61', 'C67', 'C73'\}} \\
\texttt{ARGUMENT: \{'Predlagam nabor, ki izkoristi 5510 od 5514 na voljo napora. To je visoka izkoritvenost virov, hkrati pa je nabor raznolik in zajema predmete iz vseh kategorij. Menim, da je to solidna osnova za dogovor.'\}} \\
            }
        }
    }
    & \\ \\

    \theutterance \stepcounter{utterance}  
    & & & \multicolumn{4}{p{0.6\linewidth}}{
        \cellcolor[rgb]{0.9,0.9,0.9}{
            \makecell[{{p{\linewidth}}}]{
                \texttt{\tiny{[GM$\langle$P2]}}
                \texttt{ARGUMENT: \{"Predlagani nabor sicer dobro izkorisca limit, a iz moje perspektive vsebuje prevec predmetov z nizko vrednostjo na enoto napora (npr. C42, C67, C10). Tudi C39 in C16 nista najbolj ucinkovita. Zamenjava teh predmetov z manjsimi, a vrednejimi (npr. C70, B40, A43, B42) bi dala bistveno vecjo skupno pomembnost pri podobnem ali manjsem naporu. To je za mene sprejemljiv predlog le, ce bi bili vkljuceni vsaj nekateri od kljucnih predmetov."\}} \\
            }
        }
    }
     \\ \\

    \theutterance \stepcounter{utterance}  
    & & \multicolumn{4}{p{0.6\linewidth}}{
        \cellcolor[rgb]{0.9,0.9,0.9}{
            \makecell[{{p{\linewidth}}}]{
                \texttt{\tiny{[P1$\langle$GM]}}
                \texttt{ARGUMENT: \{"Predlagani nabor sicer dobro izkorisca limit, a iz moje perspektive vsebuje prevec predmetov z nizko vrednostjo na enoto napora (npr. C42, C67, C10). Tudi C39 in C16 nista najbolj ucinkovita. Zamenjava teh predmetov z manjsimi, a vrednejimi (npr. C70, B40, A43, B42) bi dala bistveno vecjo skupno pomembnost pri podobnem ali manjsem naporu. To je za mene sprejemljiv predlog le, ce bi bili vkljuceni vsaj nekateri od kljucnih predmetov."\}} \\
            }
        }
    }
    & \\ \\

    \theutterance \stepcounter{utterance}  
    & \multicolumn{4}{p{0.6\linewidth}}{
        \cellcolor[rgb]{0.9,0.9,0.9}{
            \makecell[{{p{\linewidth}}}]{
                \texttt{\tiny{[P1$\rangle$GM]}}
                \texttt{PREDLOG: \{'A21', 'A67', 'B36', 'C16', 'C39', 'C42', 'C61', 'C67', 'C70', 'C73'\}} \\
\texttt{ARGUMENT: \{'Sprejel sem vašo kritiko o nizki učinkovitosti manjših predmetov in zamenjal C10 (183 napor, 279 vrednost) z C70 (732 napor, 828 vrednost). C70 ima boljšo vrednost na enoto napora (1,13 vs 1,52) in poveča skupno pomembnost nabora, hkrati pa ostajamo znotraj limita (skupni napor: 5117). Ohranil sem raznolikost in visoko izkoritvenost virov.'\}} \\
            }
        }
    }
    & & \\ \\

    \theutterance \stepcounter{utterance}  
    & & \multicolumn{4}{p{0.6\linewidth}}{
        \cellcolor[rgb]{0.9,0.9,0.9}{
            \makecell[{{p{\linewidth}}}]{
                \texttt{\tiny{[GM$\rangle$P2]}}
                \texttt{PREDLOG: \{'A21', 'A67', 'B36', 'C16', 'C39', 'C42', 'C61', 'C67', 'C70', 'C73'\}} \\
\texttt{ARGUMENT: \{'Sprejel sem vašo kritiko o nizki učinkovitosti manjših predmetov in zamenjal C10 (183 napor, 279 vrednost) z C70 (732 napor, 828 vrednost). C70 ima boljšo vrednost na enoto napora (1,13 vs 1,52) in poveča skupno pomembnost nabora, hkrati pa ostajamo znotraj limita (skupni napor: 5117). Ohranil sem raznolikost in visoko izkoritvenost virov.'\}} \\
            }
        }
    }
    & \\ \\

    \theutterance \stepcounter{utterance}  
    & & & \multicolumn{4}{p{0.6\linewidth}}{
        \cellcolor[rgb]{0.9,0.9,0.9}{
            \makecell[{{p{\linewidth}}}]{
                \texttt{\tiny{[GM$\langle$P2]}}
                \texttt{ARGUMENT: \{"C70 je izboljsava, a C73 (211/307=1,45), C67 (162/258=1,59), C42 (50/146=2,92) in C16 (259/355=1,37) ostajajo slabi izbor za moje cilje. Prav tako C39 (684/780=1,14) in A67 (567/663=1,17) nista konkurencnosposobna. Predlagam zamenjavo tega bloka (skupaj \textasciitilde{}1833 napor, \textasciitilde{}2349 vrednost) z A43 (929/1025=1,10) in B42 (861/957=1,11) {-} skupaj 1790 napor, 1982 vrednost. Ne glede na malo nizjic ratio, ta dva prineseta absolutno vecjo vrednost in oslobita \textasciitilde{}43 napor za dodaten majhni predmet (npr. B36). To bi bilo za mene minimalni pogoj za napredek."\}} \\
            }
        }
    }
     \\ \\

    \theutterance \stepcounter{utterance}  
    & & \multicolumn{4}{p{0.6\linewidth}}{
        \cellcolor[rgb]{0.9,0.9,0.9}{
            \makecell[{{p{\linewidth}}}]{
                \texttt{\tiny{[P1$\langle$GM]}}
                \texttt{ARGUMENT: \{"C70 je izboljsava, a C73 (211/307=1,45), C67 (162/258=1,59), C42 (50/146=2,92) in C16 (259/355=1,37) ostajajo slabi izbor za moje cilje. Prav tako C39 (684/780=1,14) in A67 (567/663=1,17) nista konkurencnosposobna. Predlagam zamenjavo tega bloka (skupaj \textasciitilde{}1833 napor, \textasciitilde{}2349 vrednost) z A43 (929/1025=1,10) in B42 (861/957=1,11) {-} skupaj 1790 napor, 1982 vrednost. Ne glede na malo nizjic ratio, ta dva prineseta absolutno vecjo vrednost in oslobita \textasciitilde{}43 napor za dodaten majhni predmet (npr. B36). To bi bilo za mene minimalni pogoj za napredek."\}} \\
            }
        }
    }
    & \\ \\

    \theutterance \stepcounter{utterance}  
    & \multicolumn{4}{p{0.6\linewidth}}{
        \cellcolor[rgb]{0.9,0.9,0.9}{
            \makecell[{{p{\linewidth}}}]{
                \texttt{\tiny{[P1$\rangle$GM]}}
                \texttt{PREDLOG: \{'A21', 'A43', 'A67', 'B36', 'B42', 'C42', 'C61', 'C67', 'C70'\}} \\
\texttt{ARGUMENT: \{'Sprejel sem vašo predlog za zamenjavo: odpravil sem C16, C39, C73 in C73 (skupaj 1833 napor, 2349 vrednost) in jih zamenjal z A43 in B42 (skupaj 1790 napor, 1982 vrednost). Hkrati sem ohranil C70 in dodal B36, kar dvigne skupno pomembnost nad prejšnji predlog. Skupni napor je 5455, kar je blizu limitu. Menim, da je to bilančno izboljšava, ki upošteva vaše zahteve po višji absolutni vrednosti.'\}} \\
            }
        }
    }
    & & \\ \\

    \theutterance \stepcounter{utterance}  
    & & \multicolumn{4}{p{0.6\linewidth}}{
        \cellcolor[rgb]{0.9,0.9,0.9}{
            \makecell[{{p{\linewidth}}}]{
                \texttt{\tiny{[GM$\rangle$P2]}}
                \texttt{PREDLOG: \{'A21', 'A43', 'A67', 'B36', 'B42', 'C42', 'C61', 'C67', 'C70'\}} \\
\texttt{ARGUMENT: \{'Sprejel sem vašo predlog za zamenjavo: odpravil sem C16, C39, C73 in C73 (skupaj 1833 napor, 2349 vrednost) in jih zamenjal z A43 in B42 (skupaj 1790 napor, 1982 vrednost). Hkrati sem ohranil C70 in dodal B36, kar dvigne skupno pomembnost nad prejšnji predlog. Skupni napor je 5455, kar je blizu limitu. Menim, da je to bilančno izboljšava, ki upošteva vaše zahteve po višji absolutni vrednosti.'\}} \\
            }
        }
    }
    & \\ \\

    \theutterance \stepcounter{utterance}  
    & & & \multicolumn{4}{p{0.6\linewidth}}{
        \cellcolor[rgb]{0.9,0.9,0.9}{
            \makecell[{{p{\linewidth}}}]{
                \texttt{\tiny{[GM$\langle$P2]}}
                \texttt{ARGUMENT: \{"Vas predlog je korak naprej (A43 in B42 sta odlicna), a vsebuje vedno tri predmete, ki za mene imajo unizko vrednost na enoto napora: A67 (1,17), C67 (1,59) in C42 (2,92). C42 je edini z visoko učinkovitostjo, a prispevek k absolutni vrednosti je zanemarljiv (146). A67 in C67 skupaj porabita 729 napor za le 921 vrednost. Predlagam zamenjavo A67, C67 in C42 (skupaj 779 napor, 1067 vrednost) z A62 (967 napor, 1063 vrednost). Napor naraste za 188 (novo stanje 5643 {-} PRECORAK), zato bi morali odpraviti se B36 (307 napor, 403 vrednost). Koncni nabor: A21, A43, A62, B42, C61, C70. Napor: 5147. To je moja mejna ponudba {-} majhna, a visoko vredna skupina."\}} \\
            }
        }
    }
     \\ \\

    \theutterance \stepcounter{utterance}  
    & & \multicolumn{4}{p{0.6\linewidth}}{
        \cellcolor[rgb]{0.9,0.9,0.9}{
            \makecell[{{p{\linewidth}}}]{
                \texttt{\tiny{[P1$\langle$GM]}}
                \texttt{ARGUMENT: \{"Vas predlog je korak naprej (A43 in B42 sta odlicna), a vsebuje vedno tri predmete, ki za mene imajo unizko vrednost na enoto napora: A67 (1,17), C67 (1,59) in C42 (2,92). C42 je edini z visoko učinkovitostjo, a prispevek k absolutni vrednosti je zanemarljiv (146). A67 in C67 skupaj porabita 729 napor za le 921 vrednost. Predlagam zamenjavo A67, C67 in C42 (skupaj 779 napor, 1067 vrednost) z A62 (967 napor, 1063 vrednost). Napor naraste za 188 (novo stanje 5643 {-} PRECORAK), zato bi morali odpraviti se B36 (307 napor, 403 vrednost). Koncni nabor: A21, A43, A62, B42, C61, C70. Napor: 5147. To je moja mejna ponudba {-} majhna, a visoko vredna skupina."\}} \\
            }
        }
    }
    & \\ \\

    \theutterance \stepcounter{utterance}  
    & \multicolumn{4}{p{0.6\linewidth}}{
        \cellcolor[rgb]{0.9,0.9,0.9}{
            \makecell[{{p{\linewidth}}}]{
                \texttt{\tiny{[P1$\rangle$GM]}}
                \texttt{PREDLOG: \{'A21', 'A43', 'A62', 'B42', 'C61', 'C70'\}} \\
\texttt{ARGUMENT: \{'Sprejel sem vašo predlog za zamenjavo. Odstranil sem A67, C67, C42 in B36 (skupaj 1086 napor, 1470 vrednost) in jih zamenjal z A62 (967 napor, 1063 vrednost). Kreirani nabor sestavljajo samo predmeti z visokim absolutnim prispevkom: A21, A43, A62, B42, C61, C70. Skupni napor znaša 5147, kar je varno znotraj limita. To je kompaktna, visoko vredna skupina, ki izkoristi vira učinkovito. Mislim, da je to odlična osnova za končni dogovor.'\}} \\
            }
        }
    }
    & & \\ \\

    \theutterance \stepcounter{utterance}  
    & & \multicolumn{4}{p{0.6\linewidth}}{
        \cellcolor[rgb]{0.9,0.9,0.9}{
            \makecell[{{p{\linewidth}}}]{
                \texttt{\tiny{[GM$\rangle$P2]}}
                \texttt{PREDLOG: \{'A21', 'A43', 'A62', 'B42', 'C61', 'C70'\}} \\
\texttt{ARGUMENT: \{'Sprejel sem vašo predlog za zamenjavo. Odstranil sem A67, C67, C42 in B36 (skupaj 1086 napor, 1470 vrednost) in jih zamenjal z A62 (967 napor, 1063 vrednost). Kreirani nabor sestavljajo samo predmeti z visokim absolutnim prispevkom: A21, A43, A62, B42, C61, C70. Skupni napor znaša 5147, kar je varno znotraj limita. To je kompaktna, visoko vredna skupina, ki izkoristi vira učinkovito. Mislim, da je to odlična osnova za končni dogovor.'\}} \\
            }
        }
    }
    & \\ \\

    \theutterance \stepcounter{utterance}  
    & & & \multicolumn{4}{p{0.6\linewidth}}{
        \cellcolor[rgb]{0.9,0.9,0.9}{
            \makecell[{{p{\linewidth}}}]{
                \texttt{\tiny{[GM$\langle$P2]}}
                \texttt{STRINJAM SE: \{'A21', 'A43', 'A62', 'B42', 'C61', 'C70'\}} \\
            }
        }
    }
     \\ \\

    \theutterance \stepcounter{utterance}  
    & & & \multicolumn{2}{p{0.3\linewidth}}{
        \cellcolor[rgb]{0.9,0.9,0.9}{
            \makecell[{{p{\linewidth}}}]{
                \texttt{\tiny{[GM$|$GM]}}
                \texttt{game successful} \\
            }
        }
    }
    & & \\ \\

    \theutterance \stepcounter{utterance}  
    & & & \multicolumn{2}{p{0.3\linewidth}}{
        \cellcolor[rgb]{0.9,0.9,0.9}{
            \makecell[{{p{\linewidth}}}]{
                \texttt{\tiny{[GM$|$GM]}}
                \texttt{end game} \\
            }
        }
    }
    & & \\ \\

\end{supertabular}
}

\end{document}
